\section{Conclusion}

Predictive learning and planning impose different constraints on a representation. Controlled-future equivalence identifies which language histories may share a state. Transition prediction preserves how actions move between those states. Neither condition fixes a goal-directed metric. GAR adds the missing local order supervision by comparing counterfactual action consequences.

The experiments support this mechanism-level account. A non-generative JEPA is a viable conditional planner, although the strongest sentence model performs better. Latent prediction alone is substantially weaker than predictive GAR. Counterfactual successor prediction and counterfactual ordering are complementary. Raw distance is weaker than the learned GAR energy, while a matched direct GAR scorer remains competitive in the current in-distribution test.

The final point is as important as the positive result. The evidence supports GAR supervision, but does not yet show that successor factorization is uniquely valuable. Controlled semantic interventions, local geometry correlations, matched generalization tests, a non-arithmetic domain, and a valid test-time scaling protocol determine whether the stronger representation claim survives.
